\section{Introduction}\label{sec:introduction}

The equilibrium spin autocorrelation measures how much a system remembers
its initial configuration after averaging over its stationary dynamics.
At the critical temperature of the Sherrington--Kirkpatrick model, this
memory lives on a time scale that diverges with the number of spins.
We determine its limiting random function and prove that Gaussian and
Rademacher interactions give the same law on the stated clock.

\subsection{Model and result}
Let $D$ be either the standard Gaussian distribution or the uniform
distribution on $\{-1,1\}$. For $n\ge2$, let $(J_{ij})_{i<j}$ be independent
with law $D$, put $J_{ji}=J_{ij}$ and $J_{ii}=0$, and define
\begin{equation}\label{eq:intro-gibbs}
 \pi_{n,J}(\sigma)=\frac1{Z_{n,J}}
 \exp\left(\frac1{\sqrt n}\sum_{i<j}J_{ij}\sigma_i\sigma_j\right),
 \qquad \sigma\in\{-1,1\}^n.
\end{equation}
This is the zero-field model at inverse temperature one. Each site has an
independent rate-one clock. When the clock at site $i$ rings, its spin is
resampled from its conditional law under $\pi_{n,J}$:
\[
 \P\{\sigma_i\gets a\mid\sigma_{-i}\}
 =\frac{\exp(a h_i(\sigma))}{2\cosh h_i(\sigma)},
 \qquad h_i(\sigma)=\frac1{\sqrt n}\sum_{j\ne i}J_{ij}\sigma_j.
\]
One unit of this clock contains $n$ site updates on average.
Write $P_t^J$ for this reversible semigroup. Its stationary spin
autocorrelation and its critical rescaling are
\begin{align}
 C_{n,J}(t)&=\frac1n\sum_{i=1}^n
    \langle\sigma_i,P_t^J\sigma_i\rangle_{\pi_{n,J}},
       \label{eq:intro-correlation}\\
 A_{n,J}(s)&=n^{1/3}C_{n,J}(s n^{2/3}),\qquad s>0.
       \label{eq:intro-rescaling}
\end{align}
Thermal fluctuations and the randomness of the dynamics have already been
averaged in these quantities. Their only remaining randomness is the
interaction matrix.

Let $C_{\mathrm{loc}}((0,\infty))$ denote the continuous real functions,
with uniform convergence on every compact subinterval of $(0,\infty)$.

\begin{theorem}[Functional universality]\label{thm:main}
There is a Borel probability law $Q_{\mathrm{crit}}$ on
$C_{\mathrm{loc}}((0,\infty))$ such that, for each of the two choices of $D$,
\[
 \Law_D(A_{n,J})\Longrightarrow Q_{\mathrm{crit}}
 \qquad(n\to\infty).
\]
The law $Q_{\mathrm{crit}}$ is not a point mass. It assigns probability one
to nonzero, nonnegative, nonincreasing functions $f$ satisfying
$\lim_{s\to\infty}f(s)=0$. The Gaussian and Rademacher models use exactly
the normalization in \eqref{eq:intro-rescaling}.
\end{theorem}

The restriction to positive times is necessary: $A_{n,J}(0)=n^{1/3}$.
At every fixed positive time, the limiting value has unbounded support,
as proved in Proposition~\ref{prop:functional-randomness}.
The limiting law is described explicitly in Section~\ref{sec:edge-law}.
It is the law of a sum of coordinate autocorrelations for a reversible semigroup in
countably many edge coordinates.
The equilibrium measure of those coordinates is Gaussian conditioned on a
renormalized quadratic constraint; its parameters are the random GOE edge
points. A single positive deterministic coefficient accounts for the
relaxation of microscopic spin gradients. Section~\ref{sec:microscopic}
characterizes that coefficient by a determinate moment problem.

\subsection{History and relation to earlier work}
Sherrington and Kirkpatrick introduced their infinite-range spin-glass
model in 1975~\cite{SK1975}. Dynamical mean-field theory subsequently studied
relaxation through correlation and response, notably in the soft-spin
framework of Sompolinsky and Zippelius~\cite{SompolinskyZippelius1982}.
For critical Ising heat-bath dynamics, Billoire and
Campbell~\cite[Equation~(7)]{BilloireCampbell2011} tested numerically the
finite-size scaling form
\[
 q_c(n,t)\sim n^{-1/3}F(t/n^{2/3}),
\]
with time measured in updates per spin. This is the direct numerical
precedent for \eqref{eq:intro-rescaling}. A disorder-averaged scaling curve
does not determine whether individual sample curves concentrate; our result
identifies their limiting probability law.

Spherical models made the role of the spectral edge especially clear.
Kosterlitz, Thouless, and Jones~\cite{KosterlitzThoulessJones1976} identified
the edge mechanism in the spherical spin-glass transition.
Cugliandolo and Dean~\cite{CugliandoloDean1995} analyzed Langevin relaxation
and aging, and Ben Arous, Dembo, and Guionnet~\cite{BenArousDemboGuionnet2001}
proved limiting dynamics and aging results for a soft spherical constraint.
At finite size, Fyodorov, Perret, and Schehr~\cite{FyodorovPerretSchehr2015}
related non-self-averaging on the $n^{2/3}$ time scale to edge fluctuations
in zero-temperature spherical dynamics. These are spectral and dynamical
predecessors, with different temperatures, initial conditions, and spin
spaces from stationary critical Ising heat-bath dynamics.

The equilibrium comparison used here has its own history.
Comets~\cite{Comets1996} related Ising and spherical partition functions by
averaging over the eigenframe. Baik and Lee~\cite{BaikLee2016} studied
spherical free-energy fluctuations on the two sides of the transition,
and Landon~\cite{Landon2022} analyzed critical free-energy fluctuations and
the edge-scale saddle. Du and Huang~\cite{DuHuang2026FreeEnergy} obtained
a quantitative critical Ising--spherical partition comparison; their
quenched overlap theorem~\cite{DuHuang2026Overlap} describes the random
critical overlap law using Gaussian coordinates conditioned on a
centered square sum. The conditioned measure used below follows this framework.
The additional work here identifies the closed dynamical form, its
microscopic coefficient, and the functional limit across two coupling laws.

Rigorous high-temperature mixing developed through spectral functional
inequalities and localization. Relevant steps include
Bauerschmidt and Bodineau~\cite{BauerschmidtBodineau2019},
Eldan, Koehler, and Zeitouni~\cite{EldanKoehlerZeitouni2022}, and
Anari, Koehler, and Vuong~\cite{AnariKoehlerVuong2024}.
For fixed inverse temperature $\beta<1/2$, Wang~\cite{Wang2026} obtains
$O_\beta(n\log(n/\epsilon))$ random-site heat-bath mixing, uniformly in
external fields, with probability tending to one over the disorder.
At zero field, Boban, Li, and Oveis
Gharan~\cite{BobanLiOveisGharan2026} obtain a spectral gap bounded below by
a positive constant for rate-one-per-site dynamics, and polynomial mixing
for $\beta<1/2+\epsilon_0$, with a universal
$\epsilon_0>0$, again with high disorder probability.
Such fixed-subcritical estimates do not identify dynamics
at the endpoint $\beta=1$.

Two companion manuscripts supply the Gaussian estimates used below.
The critical slowing-down manuscript~\cite[Theorems~1.1 and~1.2]{CS} identifies the $n^{2/3}$
scale of covariance and linear relaxation, and gives an obstruction to
mixing from most realized equilibrium starts at smaller times. The critical
mixing manuscript~\cite[Theorem~1.1]{CM} obtains the matching mixing exponent and
functional inequalities. Its comparison of spectral caps with spherical
measures provides the starting estimates for our proof. An exponent statement
alone does not identify the dynamics at a fixed multiple of $n^{2/3}$;
we need quantitative improvements and an identification of the limiting
energy. We state the companion inputs with their precise ranges in
Section~\ref{sec:edge-law} and prove those additional steps here.

Finally, the comparison of coupling laws belongs to the Lindeberg approach
to universality~\cite{Chatterjee2006}. Carmona and Hu~\cite{CarmonaHu2006}
proved universality of the SK limiting free energy. The observable here
depends on the Gibbs law and the generator simultaneously. Its critical
derivative bounds require the quadratic-feature comparison described next.

\subsection{The proof and its main ingredients}
We first work with Gaussian couplings. Adjoining an independent Gaussian
diagonal does not affect the Gibbs law and makes the interaction matrix a
GOE matrix $W=U\Lambda U^{\mathsf T}$. The spectral coordinates
$x=n^{-1/3}U^{\mathsf T}\sigma$ isolate the fluctuations near the upper edge.
Their limiting equilibrium law is a conditioned measure $\Pi$, rather than
a product of independent Gaussian coordinates. This constraint remains
visible in the domain of the limiting gradient form.

There are two separate reasons why equilibrium convergence does not already
give the answer. First, a function of the edge coordinates can lower its
Dirichlet energy by adding a small microscopic correction. Second, convergence
of finitely many edge coordinates gives no bound on the sum over all
coordinates in \eqref{eq:intro-correlation}. The proof treats these issues
before comparing the two coupling distributions.

The starting observation is the resolvent trace. If $H$ is minus the
heat-bath generator, then for $z>0$ set
\[
 B_n(z)=\sum_{a=1}^n
   \langle x_a,(z+n^{2/3}H)^{-1}x_a\rangle.
\]
This quantity is independent of the choice of orthogonal coordinates. It
is the Stieltjes transform of the same spectral measure whose Laplace
transform is $A_{n,J}$. We establish convergence of these traces and then
recover convergence of the functions on every compact positive time interval.

The observations used below have a concrete spectral meaning. At scale
$p>0$, remove from $W$ the matrix
\[
 B_p=U\diag\bigl((\lambda_a-(2-p^2))_+\bigr)U^{\mathsf T},
 \qquad h_p=B_p\sigma+\mathsf N(0,B_p).
\]
Conditioning on $h_p$ replaces the interaction by $W-B_p$, whose eigenvalues
are capped at $2-p^2$, and adds the external field $h_p$. Decreasing $p$
retains fewer observed directions near the spectral edge.

The Gaussian argument has four stages.
\begin{enumerate}
\item Spectral observations smooth functions on the cube. We compare their
posterior means with a spherical model and control the score of the
comparison density. Sections~\ref{sec:observations} and~\ref{sec:variance}
sharpen the cap estimates to a conditional variance bound of order $p^{-2}$
at cap scale $p$, without a power loss. The gain comes from averaging the
curvature error before compressing it onto a smaller random spectral subspace.
\item In Section~\ref{sec:gradient}, slopes of these smoothings converge to
weak gradients for $\Pi$. We prove that their domain equals the closure of
smooth cylinder functions. This identifies the closed energy form on the
conditioned equilibrium measure.
\item Section~\ref{sec:microscopic} identifies the energy remaining after
microscopic correction. Fixed-order contractions specify a deterministic
positive measure $\omega$ on $[0,\infty)$. Polynomial approximation isolates
its atom $c_*=\omega(\{0\})$, producing matching lower bounds and recovery
sequences for the limiting form $c_*\int|\nabla F|^2\,d\Pi$.
\item Section~\ref{sec:trace} bounds the tail of the resolvent trace by vector
integration by parts. The summable decay of the resulting coefficients makes
finite-coordinate convergence sufficient for the whole observable.
\end{enumerate}

For Rademacher couplings we compare basis-independent traces directly.
The Lindeberg argument in Section~\ref{sec:comparison} first transfers a
trace built from quadratic spin features. Its vanishing controls the
fourth-order remainders for the spin trace. This order of comparison is what
allows Gaussian and sign couplings, whose fourth moments differ, to have
the same limit. Finally, Section~\ref{sec:functional} proves the functional
convergence, decay, and surviving disorder randomness. Appendix~\ref{sec:support}
explains why almost-sure decay is compatible with nondecaying functions in
the closed topological support.
